% Part:sets-functions-relations
% Chapter: size-of-sets
% Section: comparing-sizes

\documentclass[../../../include/open-logic-section]{subfiles}

\begin{document}

\olfileid{sfr}{siz}{car}

\olsection{وکھرے سائز دے سیٹ تے کانتور دا قضیہ}

\begin{explain}
اسی ایس خیال نوں ٹھیک طور تے بیان کر دتا اے کہ دو سیٹاں دا سائز اکّو
اے۔ اسی ایس خیال نوں وی ٹھیک طور تے بیان کر سکدے آں کہ اک سیٹ دوجے
توں چھوٹا اے۔ ”توں چھوٹا (یا ہم تعداد) اے“ دی ساڈی تعریف سیٹاں دے
وچکار !!a{bijection} دی بجائے پہلے سیٹ توں دوجے ول !!a{injection}
دی منگ کرے گی۔ جے ایہو جہا فنکشن موجود ہووے، تاں پہلے سیٹ دا سائز
دوجے سیٹ دے سائز توں گھٹ یا اوہدے برابر اے۔ بدیہی طور تے، اک سیٹ
توں دوجے سیٹ ول !!a{injection} ایہ ضمانت دیندا اے کہ فنکشن دی حاصل
قدراں وچ گھٹ توں گھٹ اوہنے !!{element}s نیں جِنّے دائرۂ تعریف وچ،
کیوں جے دائرۂ تعریف دے کوئی دو !!{element}s حاصل قدراں دے اکو
!!{element} ول نہیں بھیجے جاندے۔
\end{explain}

\begin{defn}
$A$،~$B$ توں \emph{وڈا نہیں}، جیہنوں $\cardle{A}{B}$ لکھدے آں،
عین اوس ویلے اے جدوں کوئی !!a{injection} $f \colon A \to B$ موجود ہووے۔
\end{defn}

ایہ صاف اے کہ ایہ اک انعکاسی تے متعدی تعلق اے، پر متناظر نہیں (ایہ
گل مشق دے طور تے چھڈی گئی اے)۔ اسی اک ایہو جہا تصور وی متعارف کرا
سکدے آں جیہڑا دسدا اے کہ اک سیٹ دوجے توں (سختی نال) چھوٹا اے۔

\begin{defn}
$A$،~$B$ توں \emph{چھوٹا} اے، جیہنوں $\cardless{A}{B}$ لکھدے آں،
عین اوس ویلے جدوں کوئی !!a{injection}~$f\colon A \to B$ موجود ہووے
پر کوئی !!{bijection}~$g\colon A \to B$ موجود نہ ہووے، یعنی
$\cardle{A}{B}$ تے $\cardneq{A}{B}$۔
\end{defn}

ایہ صاف اے کہ ایہ تعلق غیر انعکاسی تے متعدی اے۔ (ایہ گل مشق دے طور
تے چھڈی گئی اے۔) ایس علامت نال اسی آکھ سکدے آں کہ سیٹ~$A$،
!!{enumerable} عین اوس ویلے اے جدوں $\cardle{A}{\Nat}$، تے~$A$،
!!{nonenumerable} عین اوس ویلے اے جدوں $\cardless{\Nat}{A}$۔ ایس نال
اسی
\oliflabeldef{sfr:siz:nen-alt:thm:nonenum-pownat}{%
\olref[sfr][siz][nen-alt]{thm:nonenum-pownat}
نوں ایس مشاہدے دے طور تے فیر بیان کر سکدے آں کہ
$\cardless{\Nat}{\Pow{\Nat}}$}{%
\olref[sfr][siz][nen]{thm:nonenum-pownat}
نوں ایس مشاہدے دے طور تے فیر بیان کر سکدے آں کہ
$\cardless{\PosInt}{\Pow{\PosInt}}$}۔ اصل وچ، \citet{Cantor1892} نے
ثابت کیتا کہ ایہ پچھلی گل \emph{پوری طرح عام} اے:

\begin{thm}[کانتور]\ollabel{thm:cantor}
$\cardless{A}{\Pow{A}}$، ہر سیٹ~$A$ لئی۔
\end{thm}

\begin{proof}
نقشہ $f(x) = \{x\}$ اک !!a{injection} $f \colon A \to \Pow{A}$ اے،
کیوں جے جے $x \neq y$، تاں توسیعیت موجب $\{x\} \neq \{y\}$ وی اے،
تے ایس لئی $f(x) \neq f(y)$۔ سو ساڈے کول
$\cardle{A}{\Pow{A}}$ اے۔

\begin{editorial}
جے \olref[nen]{sec} موجود ہووے تاں اسی ہولا ثبوت پیش کردے آں، نہیں
تاں \olref[nen-alt]{sec} نال میل کھان والا مختصر ثبوت۔
\end{editorial}

\oliflabeldef{sfr:siz:nen:sec}{%
  ہُن اسی وکھاواں گے کہ کوئی !!a{surjective} فنکشن~$g\colon A \to
  \Pow{A}$ موجود نہیں ہو سکدا، !!a{bijective} ہونا تاں دور دی گل اے،
  تے ایس لئی $\cardneq{A}{\Pow{A}}$۔ فرض کرو کہ
  $g\colon A \to \Pow{A}$ اے۔ کیوں جے $g$ کُل اے، ہر $x \in A$ نوں
  اک ذیلی سیٹ $g(x) \subseteq A$ ول بھیجیا جاندا اے۔ اسی وکھا سکدے آں
  کہ $g$ ہر ہدف قدر لین والا نہیں ہو سکدا۔ ایہ کرن لئی اسی اک ذیلی سیٹ
  ~$\overline{A} \subseteq A$ تعریف کردے آں جیہڑا اپنی تعریف موجب
  ~$g$ دی حاصل قدراں وچ نہیں ہو سکدا۔ رکھو
  \[
  \overline{A} = \Setabs{x \in A}{x \notin g(x)}.
  \]
  کیوں جے $g(x)$ ہر $x \in A$ لئی تعریف اے، $\overline{A}$ صاف
  طور تے~$A$ دا اک درست تعریف ہویا ذیلی سیٹ اے۔ پر ایہ~$g$ دی حاصل
  قدراں وچ نہیں ہو سکدا۔ کوئی من مانا $x \in A$ لو؛ اسی وکھاواں گے
  کہ $\overline{A} \neq g(x)$۔ جے $x \in g(x)$، تاں اوہ
  $x \notin g(x)$ نوں پورا نہیں کردا، تے ایس لئی~$\overline{A}$ دی
  تعریف موجب $x \notin \overline{A}$۔ جے $x \in \overline{A}$، تاں
  اوہنوں~$\overline{A}$ دی تعریفی خاصیت، یعنی $x \in A$ تے
  $x \notin g(x)$، پوری کرنی پئے گی۔ کیوں جے $x$ من مانا سی، ایہ
  وکھاؤندا اے کہ ہر $x \in A$ لئی، $x \in g(x)$ عین اوس ویلے اے
  جدوں $x \notin \overline{A}$، تے ایس لئی
  $g(x) \neq \overline{A}$۔ ہور لفظاں وچ، $\overline{A}$،~$g$ دی
  حاصل قدراں وچ نہیں ہو سکدا، جیہڑا~$g$ دے ہر ہدف قدر لین والے ہون
  دے مفروضے نال تضاد اے۔
  % OLSIZ-007 ماخذ دی درستی: من مانا x، سیٹ A وچوں چُنیا گیا اے؛ نتیجہ A دے ہر x لئی اے، نہ کہ صرف خط والے A دے رکناں لئی۔
}{ایہ وکھاؤنا باقی اے کہ $\cardneq{A}{\Pow{A}}$۔ برہانِ خلف لئی فرض
  کرو $\cardeq{A}{\Pow{A}}$، یعنی کوئی !!{bijection}
  $g \colon A \to \Pow{A}$ موجود اے۔ ہُن غور کرو:
  \[
    D = \Setabs{x \in A}{x \notin g(x)}
  \]
  دھیان دیو کہ $D \subseteq A$، تے سو $D \in \Pow{A}$۔ کیوں جے
  $g$ اک !!a{bijection} اے، کوئی $y \in A$ موجود اے جیہدے لئی
  $g(y) = D$۔ پر ہُن:
  \[
    y \in g(y) \text{ عین اوس ویلے جدوں } y \in D
    \text{ عین اوس ویلے جدوں } y \notin g(y).
  \]
  ایہ تضاد اے؛ سو $\cardneq{A}{\Pow{A}}$۔}{}
\end{proof}

\begin{explain}
\oliflabeldef{sfr:siz:nen:thm:nonenum-pownat}{\olref{thm:cantor} دے
  ثبوت دا \olref[nen]{thm:nonenum-pownat} دے ثبوت نال موازنہ فائدہ مند
  اے۔ اوتھے اسی وکھایا سی کہ ہر فہرست $Z_1$، $Z_2$، \dots{}، جیہڑی
  ~$\PosInt$ دے ذیلی سیٹاں دی اے، اوہدے لئی عدداں دا اک
  سیٹ~$\overline{Z}$ بنایا جا
  سکدا اے جیہڑا لازماً فہرست وچ نہیں۔ فہرست وچ نہ ہون دی ضمانت ایس
  لئی سی کہ ہر $n \in \PosInt$ لئی، $n \in Z_n$ عین اوس ویلے اے
  جدوں $n \notin \overline{Z}$۔ ایس طرح ہمیشہ کوئی عدد ہوندا اے جیہڑا
  $Z_n$ تے~$\overline{Z}$ وچوں اک دا !!a{element} اے پر دوجے دا
  نہیں۔ ایتھے اسی اوسے خیال اُتے چلدے آں، پر ہُن اشاریے~$n$،~$A$
  دے !!{element}s نیں،~$\PosInt$ دی بجائے۔ سیٹ~$\overline{A}$
  نوں ایویں تعریف کیتا گیا اے کہ اوہ~$g(x)$ توں وکھرا اے، ہر
  $x \in A$ لئی، کیوں جے $x \in g(x)$ عین اوس ویلے اے جدوں
  $x \notin \overline{A}$۔ فیر ہمیشہ~$A$ دا کوئی !!a{element} ہوندا
  اے جیہڑا $g(x)$ تے~$\overline{A}$ وچوں اک دا !!a{element} اے پر
  دوجے دا نہیں۔ تے جیویں~$\overline{Z}$ فہرست $Z_1$، $Z_2$، \dots{}
  وچ نہیں ہو سکدا، اویویں~$\overline{A}$،~$g$ دی حاصل قدراں وچ نہیں
  ہو سکدا۔}{}

\oliflabeldef{sfr:siz:nen-alt:thm:nonenum-pownat}{\olref{thm:cantor}
  دے ثبوت دا \olref[nen-alt]{thm:nonenum-pownat} دے ثبوت نال موازنہ
  فائدہ مند اے۔ اوتھے اسی وکھایا سی کہ ہر فہرست $N_0$، $N_1$،
  $N_2$، \dots{}، جیہڑی~$\Nat$ دے ذیلی سیٹاں دی اے، اوہدے لئی
  عدداں دا اک سیٹ~$D$ بنایا
  جا سکدا اے جیہڑا لازماً فہرست وچ نہیں۔ فہرست وچ نہ ہون دی ضمانت
  ایس لئی سی کہ $n \in N_n$ عین اوس ویلے اے جدوں $n \notin D$، ہر
  $n \in \Nat$ لئی۔ ایتھے اسی اوسے خیال اُتے چلدے آں، پر ہُن
  اشاریے~$n$،~$A$ دے !!{element}s نیں،~$\Nat$ دی بجائے۔ سیٹ~$D$
  نوں ایویں تعریف کیتا گیا اے کہ اوہ~$g(x)$ توں وکھرا اے، ہر
  $x \in A$ لئی، کیوں جے $x \in g(x)$ عین اوس ویلے اے جدوں
  $x \notin D$۔}{}

ایس ثبوت دا رسل دے تناقض دے ثبوت،
\olref[sfr][set][rus]{thm:russells-paradox}، نال موازنہ وی فائدہ مند اے۔
اصل وچ کانتور دا قضیہ خود رسل دے تناقض لئی تحریک سی۔
\end{explain}

\begin{prob}
  وکھاؤ کہ کوئی !!a{injection} $g\colon \Pow{A} \to A$ موجود نہیں
  ہو سکدا، کسے وی سیٹ~$A$ لئی۔ اشارا: فرض کرو
  $g\colon \Pow{A} \to A$، !!{injective} اے۔ سیٹ
  $D = \Setabs{g(B)}{B \subseteq A \text{ تے } g(B) \notin B}$
  اُتے غور کرو۔ $x = g(D)$ رکھو۔ تضاد کڈھن لئی ایس حقیقت نوں ورتو
  کہ~$g$، !!{injective} اے۔
\end{prob}

\end{document}
